To solve this problem we must calculate gravitational potential at the center
of the cloud, letting $\phi(\infty) = 0$. For a uniform density and spherical
mass distribution we have

% FIGURE NEEDED: two concentric circles (cloud of radius R, inner sphere of radius r), with radii R, r and a short arc labelled r' marking a thin shell (2009 TH solutions booklet, p. 47)
\[
  \frac{1}{2}\,m\,v^2(r=R) - \frac{GMm}{R} = E
  = \frac{1}{2}mv^2(r=0) + \varphi(0)
\]
\hfill(4 points)
\[
  v(r=R) = 0 \qquad \& \qquad v(r=0) = \sqrt{\frac{GM}{R}}
\]
\hfill(1 point)
\[
  \varphi(0) = -\frac{1}{2}mv_0^2 - \frac{GmM}{R}
\]
\hfill(1 point)
\[
  \varphi(0) = \frac{-1}{2}\,m\left(\sqrt{\frac{GM}{R}}\right)^2
             - \frac{GmM}{R}
\]
\hfill(2 points)
\[
  \varphi(0) = \frac{-3}{2}\times\frac{GMm}{R}
\]
To escape from the cloud, the particle should have total energy equal to zero
\[
  E = \tfrac{1}{2}m{v_e}^2 + \phi(r=0)
    = \tfrac{1}{2}m{v_e}^2 - \tfrac{3}{2}\left(\frac{GMm}{R}\right) = 0
\]
\[
  {v_e}^2 = \frac{3GM}{R}
  \qquad\rightarrow\qquad
  v_e = \sqrt{\frac{3GM}{R}}
\]
\hfill(2 points)
